\section*{Chapter \ref{ch:b2:diatomic-mos} --- Molecular Orbitals of Diatomic Molecules}

\begin{solution}{exo:b2:diatomic-mos:1}
$E_+ = (-13.6 - 9.0)/1.60 = \qty{-14.1}{eV}$; $E_- = (-13.6 + 9.0)/0.40 = \qty{-11.5}{eV}$.
Two electrons in $\psi_+$: $-28.25$ against $2 \times (-13.6) = \qty{-27.2}{eV}$, bound by
\qty{1.05}{eV} in this crude model.
\end{solution}

\begin{solution}{exo:b2:diatomic-mos:2}
\ce{He2}: $(2 - 2)/2 = 0$; \ce{He2+}: $(2 - 1)/2 = 0.5$; \ce{Li2}: 1; \ce{B2}: six valence
electrons, $\sigma_{2s}^2\sigma_{2s}^{*2}\pi_{2p}^2$, \omterm{def:b2:diatomic-mos:bond-order}{bond order} 1.
\end{solution}

\begin{solution}{exo:b2:diatomic-mos:3}
\ce{B2} (two unpaired electrons in the two $\pi_{2p}$ orbitals, with s--p mixing) and
\ce{O2} (two in $\pi^*$). \ce{C2}, \ce{N2}, \ce{F2} are \omterm{def:b2:diatomic-mos:magnetism}{diamagnetic}.
\end{solution}

\begin{solution}{exo:b2:diatomic-mos:4}
Non-zero: (1s, $2\mathrm p_z$), ($2\mathrm p_x$, $2\mathrm p_x$), (2s, $2\mathrm p_z$).
Zero by symmetry: (1s, $2\mathrm p_x$) and ($2\mathrm p_x$, $2\mathrm p_y$).
\end{solution}

\begin{solution}{exo:b2:diatomic-mos:5}
Ionisation removes a bonding electron: \omterm{def:b2:diatomic-mos:bond-order}{bond order} 2.5 instead of 3, so a longer and
weaker bond, as $D_0(\ce{N2+}) = \qty{8.71}{eV} < D_0(\ce{N2}) = \qty{9.76}{eV}$ shows.
\end{solution}

\begin{solution}{exo:b2:diatomic-mos:6}
The oxygen orbitals lie lower (oxygen is more electronegative): by
\cref{thm:b2:diatomic-mos:heteronuclear}, the \omterm{def:b2:diatomic-mos:bonding}{bonding orbitals} are weighted on the
lower atom, oxygen, and the higher-lying orbitals on carbon. The highest occupied
orbital, a $\sigma$ orbital of mostly carbon character pointing away from oxygen, is
the lone pair on carbon.
\end{solution}

\begin{solution}{exo:b2:diatomic-mos:7}
Only fluorine's $2\mathrm p_z$ (along the axis) has the symmetry of hydrogen's 1s:
they form a $\sigma$ orbital weighted on F and a $\sigma^*$ weighted on H. Fluorine's
$2\mathrm p_x$, $2\mathrm p_y$ (zero overlap) and its low 2s (too far in energy)
remain nonbonding: three lone pairs. The eight valence electrons fill 2s, $\sigma$
and the two nonbonding p.
\end{solution}

\begin{solution}{exo:b2:diatomic-mos:8}
$\bar\alpha = \qty{-12}{eV}$, $\Delta = \qty{4}{eV}$, $\sqrt{4 + 9} = \qty{3.61}{eV}$:
$E = \qty{-15.6}{eV}$ and \qty{-8.4}{eV}. For the lower level $(\alpha_B - E)c_B = -\beta
c_A$: $1.61c_B = 3.0c_A$, $c_A/c_B = 0.535$, $c_B^2 = 1/(1 + 0.286) = 0.78$.
\end{solution}

\begin{solution}{exo:b2:diatomic-mos:9}
Eight valence electrons: $\sigma_{2s}^2\sigma_{2s}^{*2}\pi_{2p}^4$, the $\sigma_{2p}$ orbital
(pushed above $\pi_{2p}$ by mixing) being empty. \omterm{def:b2:diatomic-mos:bond-order}{Bond order} $(6 - 2)/2 = 2$, from the
two filled $\pi$ orbitals.
\end{solution}

\begin{solution}{exo:b2:diatomic-mos:10}
NO, eleven valence electrons, one in $\pi^*$: \omterm{def:b2:diatomic-mos:bond-order}{bond order} 2.5. \ce{NO+}, ten: \omterm{def:b2:diatomic-mos:bond-order}{bond order}
3. $D_0(\ce{NO+}) = 6.51 + 13.62 - 9.26 = \qty{10.87}{eV}$: much stronger than NO, since
the electron removed was antibonding.
\end{solution}

\begin{solution}{exo:b2:diatomic-mos:11}
$E_+ + E_- = \dfrac{(\alpha + \beta)(1 - S) + (\alpha - \beta)(1 + S)}{1 - S^2} = \dfrac{2(\alpha -
\beta S)}{1 - S^2}$. It exceeds $2\alpha$ when $\alpha - \beta S > \alpha - \alpha S^2$, that is
$\beta < \alpha S$ (dividing by $-S < 0$). With four electrons, two in each orbital, the
total energy is above that of the separated orbitals: two filled orbitals repel.
\end{solution}

\begin{solution}{exo:b2:diatomic-mos:12}
\omterm{def:b2:diatomic-mos:bond-order}{Bond orders} 2.5, 2, 1.5, 1: in this order the bonds lengthen and weaken. Hydrogen
peroxide contains the peroxide O--O single bond (\ce{O2^2-} unit); \ce{KO2} contains
the superoxide ion \ce{O2-}, \omterm{def:b2:diatomic-mos:magnetism}{paramagnetic}.
\end{solution}

\begin{solution}{pb:b2:diatomic-mos:1}
\textbf{1.} Twelve.
\textbf{2.} $\sigma_{2s}^2\,\sigma_{2s}^{*2}\,\sigma_{2p}^2\,\pi_{2p}^4\,\pi_{2p}^{*2}$, the two $\pi^*$
electrons in different orbitals.
\textbf{3.} $(8 - 4)/2 = 2$.
\textbf{4.} Two: \ce{O2} is \omterm{def:b2:diatomic-mos:magnetism}{paramagnetic}.
\textbf{5.} It pairs all electrons, whereas two of them sit alone in two degenerate
$\pi^*$ orbitals.
\textbf{6.} $D_0 = 2 \times 246.79 = \qty{493.6}{kJ/mol}$, \qty{5.12}{eV}.
\textbf{7.} \ce{O2+}: $\pi^{*1}$, \omterm{def:b2:diatomic-mos:bond-order}{bond order} 2.5.
\textbf{8.} \ce{O2-}: $\pi^{*3}$, \omterm{def:b2:diatomic-mos:bond-order}{bond order} 1.5.
\textbf{9.} \ce{O2^2-}: $\pi^{*4}$, \omterm{def:b2:diatomic-mos:bond-order}{bond order} 1.
\textbf{10.} \ce{O2+} one, \ce{O2-} one, \ce{O2^2-} none.
\textbf{11.} \ce{O2+} $<$ \ce{O2} $<$ \ce{O2-} $<$ \ce{O2^2-}.
\textbf{12.} The peroxide ion.
\textbf{13.} \ce{O2+}, \ce{O2} and \ce{O2-}.
\textbf{14.} From a $\pi^*$ orbital.
\textbf{15.} The $\pi^*$ orbital, antibonding, lies above the 2p level of the atom: its
electron is less bound.
\textbf{16.} Two paths from \ce{O2} to $\ce{O} + \ce{O+} + \mathrm e^-$: dissociate then
ionise an atom, $D_0(\ce{O2}) + I(\ce{O})$; or ionise the molecule then dissociate the ion,
$I(\ce{O2}) + D_0(\ce{O2+})$. Equal: $D_0(\ce{O2+}) = 5.12 + 13.62 - 12.07 = \qty{6.66}{eV}$.
\textbf{17.} Larger than $D_0(\ce{O2}) = \qty{5.12}{eV}$: removing an antibonding electron
strengthens the bond (order 2.5).
\textbf{18.} $D_0(\ce{N2}) = 2 \times 470.82/96.485 = \qty{9.76}{eV}$; $D_0(\ce{N2+}) = 9.76 +
14.53 - 15.58 = \qty{8.71}{eV}$.
\textbf{19.} The highest occupied orbital of \ce{N2} is bonding, below the 2p level of the
atom.
\textbf{20.} When the electron removed comes from an \omterm{def:b2:diatomic-mos:bonding}{antibonding orbital}.
\textbf{21.} Hund's rule (maximum number of parallel spins in degenerate orbitals).
\textbf{22.} Both $\pi^*$ electrons in the same orbital, spins paired.
\textbf{23.} Two electrons in the same orbital repel each other more, and lose the
exchange stabilisation of parallel spins.
\textbf{24.} \omterm{def:b2:diatomic-mos:bond-order}{Bond order} still 2; all electrons paired: \omterm{def:b2:diatomic-mos:magnetism}{diamagnetic}.
\textbf{25.} \ce{O2+} has \omterm{def:b2:diatomic-mos:bond-order}{bond order} $\boldsymbol{2.5}$ and a dissociation energy of
$\boldsymbol{\approx \qty{6.66}{eV}}$, against 2 and \qty{5.12}{eV} for \ce{O2}.
\end{solution}
